% theory/revision/descent.tex -- Task 5(a) (G7-4a): the proof of Theorem 5.16 written out, with no
% PARI dependence.  Verified step by step by check_descent.py (exact arithmetic).
%
% PLACEMENT.  Replace the proof of Theorem~\ref{thm:isolation} (manuscript ll. 1041--1046) by the
% material below.  It does not use Section 8, so it survives the split of Section 8 into a
% separate note (G7-1).  New bibliography entry (fetched from the author's page, second edition
% text with corrections; Section 3.6 "Method 1: descent using 2-isogeny", printed pp. 84--85, and
% the torsion remark on printed p. 70):
%
% @book{cremona1997,
%   author    = {Cremona, J. E.},
%   title     = {Algorithms for Modular Elliptic Curves},
%   edition   = {Second},
%   publisher = {Cambridge University Press},
%   address   = {Cambridge},
%   year      = {1997},
%   note      = {Online edition: \url{https://johncremona.github.io/book/fulltext/index.html}}
% }

\begin{proof}
\emph{The model.} Put $e_2=XY+YZ+ZX$ and let
\[
  \varphi(X:Y:Z)=\Big(-\frac{16e_2}{Z^2},\ \frac{8(X-Y)}Z\Big(-\frac{4e_2}{Z^2}-54\Big)\Big),\qquad
  \psi(x,y)=\big(25x+y:\ 25x-y:\ 4(x-216)\big).
\]
A direct substitution shows that $\varphi$ maps $C_{27/2}$ into the curve
\[
  E:\ y^2=x(x+9)(x+384)=x^3+393x^2+3456x,
\]
that $\psi$ maps $E$ into $C_{27/2}$, and that $\varphi\circ\psi$ is the identity of $E$. The
discriminant of $E$ is $16\cdot3456^2\cdot(393^2-4\cdot3456)=2^{18}3^85^6\ne0$, so $E$ is an elliptic
curve. The closure of $\psi(E)$ is an irreducible curve contained in the cubic $C_{27/2}$ and
birational to $E$, so it has geometric genus $1$; it is therefore not a line, a conic or a singular
cubic, hence it is all of $C_{27/2}$, which is a nonsingular cubic. So $\varphi$ and $\psi$ are mutually
inverse isomorphisms defined over $\Q$. As $(x,y)\to\infty$ on $E$, $y/x\to\pm\infty$ and $\psi(x,y)\to(1:-1:0)$, so
$\psi$ sends the origin of $E$ to $O=(1:-1:0)$, and $C_{27/2}(\Q)$ with base point $O$ is a group isomorphic to
$E(\Q)$. For example $\varphi(1:4:4)=(-24,360)$.

\emph{Rank.} We use the descent via $2$-isogeny in the form of
\cite[\S3.6, Method~1 and (3.6.2)]{cremona1997}. Write $E:y^2=x(x^2+cx+d)$ with $c=393$,
$d=3456=2^73^3$; the isogenous curve is $E':y^2=x(x^2+c'x+d')$ with $c'=-2c=-786$ and
$d'=c^2-4d=140625=3^25^6$. For squarefree $d_1\mid d$ let $n_1$ be the number of classes $d_1$ for
which $H(d_1):v^2=d_1u^4+cu^2+d/d_1$ has a rational point, and define $n_1'$ in the same way with
$c',d'$. Then $n_1=2^{e_1}$, $n_1'=2^{e_1'}$ and $\operatorname{rank}E(\Q)=e_1+e_1'-2$
\cite[(3.6.2)]{cremona1997}. A rational point of $H(d_1)$ gives, after clearing denominators,
integers $M,e,N$ with $N^2=d_1M^4+cM^2e^2+(d/d_1)e^4$ and $\gcd(M,e)=1$; so it suffices to exclude
such solutions.
\begin{enumerate}[label=(\roman*)]
\item On $E$, $d_1\in\{\pm1,\pm2,\pm3,\pm6\}$. For $d_1\in\{\pm2,\pm3\}$ we have $d_1\equiv\pm2\pmod5$, a quadratic non-residue,
  $d/d_1\equiv d_1^{-1}\pmod5$ since $d\equiv1$, so $d_1+d_1^{-1}\equiv0$, and $c\equiv3\pmod5$. Since $M^4,e^4\in\{0,1\}$ modulo $5$,
  with $5$ dividing at most one of $M,e$, the right side is $\equiv d_1$, $\equiv d_1^{-1}$, or
  $\equiv d_1+3M^2e^2+d_1^{-1}\equiv3M^2e^2\pmod5$; each is a quadratic non-residue modulo $5$
  ($M^2e^2\equiv\pm1$), so there is no solution. The classes $\pm1,\pm6$ are attained:
  the points $(-9,0)$, $(-384,0)$, $(0,0)$ give $-9\equiv-1$, $-384\equiv-6$ and $d\equiv6$ modulo
  squares, and $O$ gives $1$. So $n_1=4$ and $e_1=2$.
\item On $E'$, $d_1\in\{\pm1,\pm3,\pm5,\pm15\}$. If $d_1<0$, all three coefficients
  $d_1,\,c'=-786,\,d'/d_1$ are negative, so the right side is negative for $(M,e)\ne(0,0)$. For
  $d_1=3,5,15$ write $f=d_1M^4-786M^2e^2+(d'/d_1)e^4$ and $f=3^vf_0$ with $3\nmid f_0$; then $v$ is odd
  or $f_0\equiv2\pmod3$, so $f$ is not a square:
  \begin{itemize}
  \item $d_1=3$: $f=3(M^4-262M^2e^2+15625e^4)$ and the bracket is $\equiv M^4-M^2e^2+e^4\equiv1\pmod3$;
  \item $d_1=5$: $f\equiv2M^4\pmod3$; if $3\mid M$, $M=3M_1$, then $f=9(45M_1^4-786M_1^2e^2+3125e^4)$ with
    bracket $\equiv2\pmod3$;
  \item $d_1=15$: $f=3g$, $g=5M^4-262M^2e^2+3125e^4\equiv2M^4-M^2e^2+2e^4\pmod3$; if $3$ divides one of
    $M,e$ then $g\equiv2\pmod3$; otherwise, with $w\equiv M^2/e^2\in\{1,4,7\}\pmod9$,
    $g\equiv e^4(5w^2-w+2)\equiv6e^4\pmod9$, so $f=9g'$ with $g'\equiv2e^4\equiv2\pmod3$.
  \end{itemize}
  So $n_1'=1$ and $e_1'=0$.
\end{enumerate}
Hence $\operatorname{rank}E(\Q)=2+0-2=0$. (The two Selmer groups are
$\{\pm1,\pm6\}$ and $\{1\}$; every class is decided by a local condition, so no
Tate--Shafarevich ambiguity arises.)

\emph{Torsion.} $E$ has good reduction at $7$ and $11$, and for an odd prime of good reduction the
reduction map on $E(\Q)_{\rm tors}$ is injective \cite[\S3.3, p.~70]{cremona1997}. Counting points,
$\#E(\mathbb F_7)=\#E(\mathbb F_{11})=12$, so $\#E(\Q)_{\rm tors}$ divides $12$. The twelve points
\[
  O,\ (0,0),\ (-9,0),\ (-384,0),\ (-24,\pm360),\ (-144,\pm2160),\ (16,\pm400),\ (216,\pm5400)
\]
lie on $E$; $(0,0)$, $(-9,0)$, $(-384,0)$ have order $2$, $(16,\pm400)$ have order $3$, and
$(-24,\pm360)$, $(-144,\pm2160)$, $(216,\pm5400)$ have order $6$ (for instance the tangent at $(16,400)$ has
slope $21$ and meets $E$ again at $(16,-400)$). So $E(\Q)=E(\Q)_{\rm tors}\cong\Z/2\times\Z/6$.

\emph{The points of $C_{27/2}$.} The twelve points $(1:0:0)$, $(0:1:0)$, $(0:0:1)$, $(1:-1:0)$,
$(0:1:-1)$, $(1:0:-1)$ and the permutations of $(1:4:4)$ and $(1:1:4)$ lie on $C_{27/2}$ and are
distinct, so they are all of $C_{27/2}(\Q)$; the positive ones are the six permutations of $(1:4:4)$
and $(1:1:4)$.

A triad with $S_1=18k$ and $R=\frac3{4k}$ has $\Lambda=\frac{27}2$, so it is a positive rational point of
$C_{27/2}$, hence a permutation of $(1:4:4)$ or $(1:1:4)$. Each has exactly one integer representative
with sum $18k$, namely $(2k,8k,8k)$ and $(3k,3k,12k)$, and both are hyperbolic since $R=\frac3{4k}<1$.
\end{proof}

% ---- Statement of Theorem 5.16: unchanged.  Remove "PARI/GP 2.17.2 [pari2172], ellrank ...
%      elltors gives torsion of order 12" from the proof; PARI may stay as a one-sentence
%      corroboration ("PARI/GP's ellrank and elltors agree") if desired.
% ---- The order-6 points of E: (-24,+-360), (-144,+-2160), (216,+-5400); order 3: (16,+-400);
%      order 2: (0,0), (-9,0), (-384,0) -- check_descent.py (d) checks each point's order.
